\documentclass[10pt,a4paper]{amsart}
\usepackage{amsmath,amssymb,amsthm,graphicx}
% Retained minimal preamble for the Stacks Project chapter more-algebra.tex (More on Algebra).
% The chapter file pinned-original-source.tex carries only \input{preamble} on its first line, so no
% preamble text lives inside the chapter itself.  The Stacks Project uses one single shared preamble for
% every chapter; the snapshot stages it identically as sources/stacks-chow/preamble.tex and
% sources/stacks-etale-cohomology/preamble.tex.  The block below is copied mechanically, line for line,
% from sources/stacks-chow/preamble.tex lines 14-18 (the Xy-pic setup) and 154-217 (the native theorem
% declarations together with the mathematical macros).  Nothing was added, removed or redefined; the
% wrapper class and the amsmath/amssymb/amsthm/graphicx base packages are supplied by the assembler.
% Chapter: more-algebra.tex -- More on Algebra.
% For commutative diagrams we use Xy-pic
\usepackage[all]{xy}

% We use 2cell for 2-commutative diagrams.
\xyoption{2cell}

% Theorem environments.
%
\theoremstyle{plain}
\newtheorem{theorem}[subsection]{Theorem}
\newtheorem{proposition}[subsection]{Proposition}
\newtheorem{lemma}[subsection]{Lemma}

\theoremstyle{definition}
\newtheorem{definition}[subsection]{Definition}
\newtheorem{example}[subsection]{Example}
\newtheorem{exercise}[subsection]{Exercise}
\newtheorem{situation}[subsection]{Situation}

\theoremstyle{remark}
\newtheorem{remark}[subsection]{Remark}
\newtheorem{remarks}[subsection]{Remarks}

\numberwithin{equation}{subsection}

% Macros
%
\def\lim{\mathop{\mathrm{lim}}\nolimits}
\def\colim{\mathop{\mathrm{colim}}\nolimits}
\def\Spec{\mathop{\mathrm{Spec}}}
\def\Hom{\mathop{\mathrm{Hom}}\nolimits}
\def\Ext{\mathop{\mathrm{Ext}}\nolimits}
\def\SheafHom{\mathop{\mathcal{H}\!\mathit{om}}\nolimits}
\def\SheafExt{\mathop{\mathcal{E}\!\mathit{xt}}\nolimits}
\def\Sch{\mathit{Sch}}
\def\Mor{\mathop{\mathrm{Mor}}\nolimits}
\def\Ob{\mathop{\mathrm{Ob}}\nolimits}
\def\Sh{\mathop{\mathit{Sh}}\nolimits}
\def\NL{\mathop{N\!L}\nolimits}
\def\CH{\mathop{\mathrm{CH}}\nolimits}
\def\proetale{{pro\text{-}\acute{e}tale}}
\def\etale{{\acute{e}tale}}
\def\QCoh{\mathit{QCoh}}
\def\Ker{\mathop{\mathrm{Ker}}}
\def\Im{\mathop{\mathrm{Im}}}
\def\Coker{\mathop{\mathrm{Coker}}}
\def\Coim{\mathop{\mathrm{Coim}}}

% Boxtimes
%
\DeclareMathSymbol{\boxtimes}{\mathbin}{AMSa}{"02}

%
% Macros for moduli stacks/spaces
%
\def\QCohstack{\mathcal{QC}\!\mathit{oh}}
\def\Cohstack{\mathcal{C}\!\mathit{oh}}
\def\Spacesstack{\mathcal{S}\!\mathit{paces}}
\def\Quotfunctor{\mathrm{Quot}}
\def\Hilbfunctor{\mathrm{Hilb}}
\def\Curvesstack{\mathcal{C}\!\mathit{urves}}
\def\Polarizedstack{\mathcal{P}\!\mathit{olarized}}
\def\Complexesstack{\mathcal{C}\!\mathit{omplexes}}
% \Pic is the operator that assigns to X its picard group, usage \Pic(X)
% \Picardstack_{X/B} denotes the Picard stack of X over B
% \Picardfunctor_{X/B} denotes the Picard functor of X over B
\def\Pic{\mathop{\mathrm{Pic}}\nolimits}
\def\Picardstack{\mathcal{P}\!\mathit{ic}}
\def\Picardfunctor{\mathrm{Pic}}
\def\Deformationcategory{\mathcal{D}\!\mathit{ef}}

% --- appended 2026-10-12 (Chief): the project own definitions of the reference/slogan/history
% display environments, copied mechanically from the shared Stacks preamble lines 8-12
% (\usepackage{verbatim} plus the three \newenvironment lines, which the chapter preamble keeps
% verbatim from the project preamble). Nothing added, removed or redefined: these environments
% carry no mathematical content and the project itself typesets them as comments. Without these
% three lines a counted statement whose own text uses them fails with Environment slogan undefined.
\usepackage{verbatim}
\newenvironment{reference}{\comment}{\endcomment}
\newenvironment{slogan}{\comment}{\endcomment}
\newenvironment{history}{\comment}{\endcomment}
\usepackage{cleveref}
\crefname{originallemma}{original lemma}{original lemmas}
\Crefname{originallemma}{Original lemma}{Original lemmas}
\makeatletter
\newlabel{more-algebra-lemma-henselization-noetherian}{{06LJ}{source}{}{sourcepointer-1}{}}
\newlabel{more-algebra-lemma-henselization-noetherian@cref}{{[originallemma][0][]06LJ}{[1][0][]source}{}{}{}}
\makeatother
\crefname{originallemma}{original lemma}{original lemmas}
\Crefname{originallemma}{Original lemma}{Original lemmas}
\makeatletter
\newlabel{lemma-dumb-properties-henselization}{{07QM}{source}{}{sourcepointer-2}{}}
\newlabel{lemma-dumb-properties-henselization@cref}{{[originallemma][0][]07QM}{[1][0][]source}{}{}{}}
\makeatother
\crefname{originallemma}{original lemma}{original lemmas}
\Crefname{originallemma}{Original lemma}{Original lemmas}
\makeatletter
\newlabel{algebra-lemma-descent-Noetherian}{{033E}{source}{}{sourcepointer-3}{}}
\newlabel{algebra-lemma-descent-Noetherian@cref}{{[originallemma][0][]033E}{[1][0][]source}{}{}{}}
\makeatother
\crefname{originallemma}{original lemma}{original lemmas}
\Crefname{originallemma}{Original lemma}{Original lemmas}
\makeatletter
\newlabel{algebra-lemma-complete-local-ring-Noetherian}{{032B}{source}{}{sourcepointer-4}{}}
\newlabel{algebra-lemma-complete-local-ring-Noetherian@cref}{{[originallemma][0][]032B}{[1][0][]source}{}{}{}}
\makeatother
\crefname{originallemma}{original lemma}{original lemmas}
\Crefname{originallemma}{Original lemma}{Original lemmas}
\makeatletter
\newlabel{algebra-lemma-completion-faithfully-flat}{{00MC}{source}{}{sourcepointer-5}{}}
\newlabel{algebra-lemma-completion-faithfully-flat@cref}{{[originallemma][0][]00MC}{[1][0][]source}{}{}{}}
\makeatother
\crefname{originallemma}{original lemma}{original lemmas}
\Crefname{originallemma}{Original lemma}{Original lemmas}
\makeatletter
\newlabel{algebra-lemma-colimit-rings-flat}{{05UU}{source}{}{sourcepointer-6}{}}
\newlabel{algebra-lemma-colimit-rings-flat@cref}{{[originallemma][0][]05UU}{[1][0][]source}{}{}{}}
\makeatother
\crefname{originallemma}{original lemma}{original lemmas}
\Crefname{originallemma}{Original lemma}{Original lemmas}
\makeatletter
\newlabel{lemma-completed-direct-sum-flat}{{06LE}{source}{}{sourcepointer-7}{}}
\newlabel{lemma-completed-direct-sum-flat@cref}{{[originallemma][0][]06LE}{[1][0][]source}{}{}{}}
\makeatother
\crefname{originallemma}{original lemma}{original lemmas}
\Crefname{originallemma}{Original lemma}{Original lemmas}
\makeatletter
\newlabel{lemma-henselization-formally-smooth}{{07QN}{source}{}{sourcepointer-8}{}}
\newlabel{lemma-henselization-formally-smooth@cref}{{[originallemma][0][]07QN}{[1][0][]source}{}{}{}}
\makeatother
\crefname{originallemma}{original lemma}{original lemmas}
\Crefname{originallemma}{Original lemma}{Original lemmas}
\makeatletter
\newlabel{lemma-formally-smooth-completion}{{07ED}{source}{}{sourcepointer-9}{}}
\newlabel{lemma-formally-smooth-completion@cref}{{[originallemma][0][]07ED}{[1][0][]source}{}{}{}}
\makeatother
\crefname{originallemma}{original lemma}{original lemmas}
\Crefname{originallemma}{Original lemma}{Original lemmas}
\makeatletter
\newlabel{algebra-lemma-sh-from-h-map}{{0C2Z}{source}{}{sourcepointer-10}{}}
\newlabel{algebra-lemma-sh-from-h-map@cref}{{[originallemma][0][]0C2Z}{[1][0][]source}{}{}{}}
\makeatother
\crefname{originallemma}{original lemma}{original lemmas}
\Crefname{originallemma}{Original lemma}{Original lemmas}
\makeatletter
\newlabel{algebra-lemma-complete-henselian}{{04GM}{source}{}{sourcepointer-11}{}}
\newlabel{algebra-lemma-complete-henselian@cref}{{[originallemma][0][]04GM}{[1][0][]source}{}{}{}}
\makeatother
\crefname{originallemma}{original lemma}{original lemmas}
\Crefname{originallemma}{Original lemma}{Original lemmas}
\makeatletter
\newlabel{algebra-lemma-strictly-henselian-functorial}{{04GU}{source}{}{sourcepointer-12}{}}
\newlabel{algebra-lemma-strictly-henselian-functorial@cref}{{[originallemma][0][]04GU}{[1][0][]source}{}{}{}}
\makeatother
\title{The henselization of a pair over a Noetherian ring is Noetherian}
\author{The Stacks Project Authors}
\date{stacks-project:more-algebra.tex:local-snapshot-20261012; GNU FDL 1.2 or later}
\begin{document}
\maketitle
\setcounter{subsection}{1}
\noindent\textit{Editorial note.} All delivered mathematical text is the author's own, copied from the pinned Stacks Project chapter file more-algebra.tex; no statement or proof step was rewritten, reordered or removed.\par\medskip
\noindent\textit{Editorial note.} The chapter file pinned-original-source.tex carries only \textbackslash{}input\{preamble\} on its first line, so no preamble text lives inside the chapter. The retained file sources/stacks-more-algebra/retained-preamble.tex is a mechanical, line-for-line copy of the project's single shared preamble as staged in the snapshot (Xy-pic setup plus theorem declarations and mathematical macros). Nothing was added, removed or redefined, and the excerpt serves only as preamble, never as a span coordinate.\par\medskip
\noindent\textit{Editorial note.} Both counted spans start at the \textbackslash{}begin line of their environment, so the statement and the author's proof are delivered whole; surrounding prose of the section is deliberately outside the counted spans and is carried as standard prerequisites instead.\par\medskip
\noindent\textit{Editorial note.} The citation fields (source\_id, source\_version, authors, source\_url, license.name, source\_title) are identical to those of every other card this lane has produced from this chapter, so the batch selector treats them as one source.\par\medskip
\noindent\textit{Editorial note.} EDITORIAL NOTE: the statement contains a comment-like 'reference' environment (the chapter defines these environments with comment/endcomment), whose content includes a citation of EGA. Because that environment is not typeset, the citation is never processed and no bibliography entry is required; the stub that had been created for it was removed, and the retained preamble is now taken in full through its current last line so that the environment definition is present. The bytes of this note are plain ASCII and contain no control characters.\par\medskip
\noindent\textit{Allowed standard prerequisites (editorial).}\begin{itemize}
\item Henselization and strict henselization of a pair, with their universal properties. defined in the same section, source lines 10600-10977, before the counted statement. Supplies the rings whose Noetherianness is asserted.
\item The strict henselization is a filtered colimit of finite etale algebras over the base. used at proof lines 11002-11030. Reduces the question to Noetherianness of finite type algebras.
\item A filtered union of finitely generated subalgebras of a ring of finite type over a Noetherian base is Noetherian. used at proof lines 11030-11060. Is the specific statement that replaces the general false statement about filtered colimits.
\end{itemize}\medskip
\ifdefined\hypertarget\hypertarget{sourcepointer-1}{}\fi
\noindent\textit{Source reference.} Original lemma 06LJ: focused-seeds/AGP/sources/stacks-snapshot-20261011/tags.txt, source lines 8468--8468. The lemma's own canonical tag, quoted so that readers can resolve the label inside the delivered text; it points at the counted statement itself and carries no proof content.\par\medskip
\ifdefined\hypertarget\hypertarget{sourcepointer-2}{}\fi
\noindent\textit{Source reference.} Original lemma 07QM: focused-seeds/AGP/sources/stacks-snapshot-20261011/tags.txt, source lines 9837--9837. Navigation pointer to a statement the counted proof refers to but does not contain; nothing load-bearing is delegated to it.\par\medskip
\ifdefined\hypertarget\hypertarget{sourcepointer-3}{}\fi
\noindent\textit{Source reference.} Original lemma 033E: focused-seeds/AGP/sources/stacks-snapshot-20261011/tags.txt, source lines 4053--4053. Navigation pointer to a statement the counted proof refers to but does not contain; nothing load-bearing is delegated to it.\par\medskip
\ifdefined\hypertarget\hypertarget{sourcepointer-4}{}\fi
\noindent\textit{Source reference.} Original lemma 032B: focused-seeds/AGP/sources/stacks-snapshot-20261011/tags.txt, source lines 4014--4014. Navigation pointer to a statement the counted proof refers to but does not contain; nothing load-bearing is delegated to it.\par\medskip
\ifdefined\hypertarget\hypertarget{sourcepointer-5}{}\fi
\noindent\textit{Source reference.} Original lemma 00MC: focused-seeds/AGP/sources/stacks-snapshot-20261011/tags.txt, source lines 818--818. Navigation pointer to a statement the counted proof refers to but does not contain; nothing load-bearing is delegated to it.\par\medskip
\ifdefined\hypertarget\hypertarget{sourcepointer-6}{}\fi
\noindent\textit{Source reference.} Original lemma 05UU: focused-seeds/AGP/sources/stacks-snapshot-20261011/tags.txt, source lines 7528--7528. Navigation pointer to a statement the counted proof refers to but does not contain; nothing load-bearing is delegated to it.\par\medskip
\ifdefined\hypertarget\hypertarget{sourcepointer-7}{}\fi
\noindent\textit{Source reference.} Original lemma 06LE: focused-seeds/AGP/sources/stacks-snapshot-20261011/tags.txt, source lines 8463--8463. Navigation pointer to a statement the counted proof refers to but does not contain; nothing load-bearing is delegated to it.\par\medskip
\ifdefined\hypertarget\hypertarget{sourcepointer-8}{}\fi
\noindent\textit{Source reference.} Original lemma 07QN: focused-seeds/AGP/sources/stacks-snapshot-20261011/tags.txt, source lines 9838--9838. Navigation pointer to a statement the counted proof refers to but does not contain; nothing load-bearing is delegated to it.\par\medskip
\ifdefined\hypertarget\hypertarget{sourcepointer-9}{}\fi
\noindent\textit{Source reference.} Original lemma 07ED: focused-seeds/AGP/sources/stacks-snapshot-20261011/tags.txt, source lines 9443--9443. Navigation pointer to a statement the counted proof refers to but does not contain; nothing load-bearing is delegated to it.\par\medskip
\ifdefined\hypertarget\hypertarget{sourcepointer-10}{}\fi
\noindent\textit{Source reference.} Original lemma 0C2Z: focused-seeds/AGP/sources/stacks-snapshot-20261011/tags.txt, source lines 15174--15174. Navigation pointer to a statement the counted proof refers to but does not contain; nothing load-bearing is delegated to it.\par\medskip
\ifdefined\hypertarget\hypertarget{sourcepointer-11}{}\fi
\noindent\textit{Source reference.} Original lemma 04GM: focused-seeds/AGP/sources/stacks-snapshot-20261011/tags.txt, source lines 5838--5838. Navigation pointer to a statement the counted proof refers to but does not contain; nothing load-bearing is delegated to it.\par\medskip
\ifdefined\hypertarget\hypertarget{sourcepointer-12}{}\fi
\noindent\textit{Source reference.} Original lemma 04GU: focused-seeds/AGP/sources/stacks-snapshot-20261011/tags.txt, source lines 5845--5845. Navigation pointer to a statement the counted proof refers to but does not contain; nothing load-bearing is delegated to it.\par\medskip
\begin{lemma}
\label{lemma-henselization-noetherian}
\begin{reference}
\cite[IV, Theorem 18.6.6 and Proposition 18.8.8]{EGA}
\end{reference}
Let $R$ be a local ring. The following are equivalent
\begin{enumerate}
\item $R$ is Noetherian,
\item $R^h$ is Noetherian, and
\item $R^{sh}$ is Noetherian.
\end{enumerate}
In this case we have
\begin{enumerate}
\item[(a)] $(R^h)^\wedge$ and $(R^{sh})^\wedge$ are Noetherian complete
local rings,
\item[(b)] $R^\wedge \to (R^h)^\wedge$ is an isomorphism,
\item[(c)] $R^h \to (R^h)^\wedge$ and $R^{sh} \to (R^{sh})^\wedge$ are flat,
\item[(d)] $R^\wedge \to (R^{sh})^\wedge$ is formally smooth in
the $\mathfrak m_{(R^{sh})^\wedge}$-adic topology,
\item[(e)] $(R^\wedge)^{sh} = R^\wedge \otimes_{R^h} R^{sh}$, and
\item[(f)] $((R^\wedge)^{sh})^\wedge = (R^{sh})^\wedge$.
\end{enumerate}
\end{lemma}

\begin{proof}
Since $R \to R^h \to R^{sh}$ are faithfully flat
(Lemma \ref{lemma-dumb-properties-henselization}),
we see that $R^h$ or $R^{sh}$ being Noetherian implies that $R$
is Noetherian, see Algebra, Lemma \ref{algebra-lemma-descent-Noetherian}.
In the rest of the proof we assume $R$ is Noetherian.

\medskip\noindent
As $\mathfrak m \subset R$ is finitely generated it follows that
$\mathfrak m^h = \mathfrak m R^h$ and $\mathfrak m^{sh} = \mathfrak mR^{sh}$
are finitely generated, see Lemma \ref{lemma-dumb-properties-henselization}.
Hence $(R^h)^\wedge$ and $(R^{sh})^\wedge$ are Noetherian by
Algebra, Lemma \ref{algebra-lemma-complete-local-ring-Noetherian}.
This proves (a).

\medskip\noindent
Note that (b) is immediate from
Lemma \ref{lemma-dumb-properties-henselization}.
In particular we see that $(R^h)^\wedge$ is flat over $R$, see
Algebra, Lemma \ref{algebra-lemma-completion-faithfully-flat}.

\medskip\noindent
Next, we show that $R^h \to (R^h)^\wedge$ is flat.
Write $R^h = \colim_i R_i$ as a directed colimit of
localizations of \'etale $R$-algebras. By
Algebra, Lemma \ref{algebra-lemma-colimit-rings-flat}
if $(R^h)^\wedge$ is flat over each $R_i$, then $R^h \to (R^h)^\wedge$ is
flat. Note that $R^h = R_i^h$ (by construction).
Hence $R_i^\wedge = (R^h)^\wedge$
by part (b) is flat over $R_i$ as desired. To finish the proof of (c)
we show that $R^{sh} \to (R^{sh})^\wedge$ is flat. To do this, by a
limit argument as above, it suffices to show that $(R^{sh})^\wedge$
is flat over $R$. Note that it follows from
Lemma \ref{lemma-dumb-properties-henselization}
that $(R^{sh})^\wedge$ is the completion of a free $R$-module.
By Lemma \ref{lemma-completed-direct-sum-flat}
we see this is flat over $R$ as desired. This finishes the proof of (c).

\medskip\noindent
At this point we know (c) is true and that $(R^h)^\wedge$ and
$(R^{sh})^\wedge$ are Noetherian. It follows from
Algebra, Lemma \ref{algebra-lemma-descent-Noetherian}
that $R^h$ and $R^{sh}$ are Noetherian.

\medskip\noindent
Part (d) follows from Lemma \ref{lemma-henselization-formally-smooth}
and Lemma \ref{lemma-formally-smooth-completion}.

\medskip\noindent
Part (e) follows from Algebra, Lemma \ref{algebra-lemma-sh-from-h-map}
and the fact that $R^\wedge$ is henselian by
Algebra, Lemma \ref{algebra-lemma-complete-henselian}.

\medskip\noindent
Proof of (f). Using (e) there is a map $R^{sh} \to (R^\wedge)^{sh}$
which induces a map $(R^{sh})^\wedge \to ((R^\wedge)^{sh})^\wedge$
upon completion. Using (e) there is a map $R^\wedge \to (R^{sh})^\wedge$.
Since $(R^{sh})^\wedge$ is strictly henselian (see above) this map
induces a map $(R^\wedge)^{sh} \to (R^{sh})^\wedge$ by
Algebra, Lemma \ref{algebra-lemma-strictly-henselian-functorial}.
Completing we obtain a map $((R^\wedge)^{sh})^\wedge \to (R^{sh})^\wedge$.
We omit the verification that these two maps are mutually inverse.
\end{proof}

\end{document}
